The Composite Planning Score (PS) gives a compact leaderboard, but it is not a replacement for the diagnostic metrics. It combines zero-shot success, retry efficiency, execution quality, grounding and precondition awareness into one bounded score so that models can be sorted quickly before inspecting their component profiles.

\begin{bmformulabox}
\begin{align*}
\mathrm{base} ={}& 0.25\cdot\mathrm{FASR} +0.20\cdot\mathrm{IWSR} +0.20\cdot\mathrm{Exec}\\
&+0.20\cdot(1-\mathrm{Hallucination}) +0.15\cdot\mathrm{PAS}
\end{align*}
\end{bmformulabox}

When a strict CoT-alignment score is available, a small optional bonus is added and the result is renormalized so PS remains in \([0,1]\).

\begin{table}[H]\small
\centering
\begin{tabularx}{\linewidth}{@{}l c Y@{}}
\toprule
\textbf{Component} & \textbf{Weight} & \textbf{Reason}\\
\midrule
FASR & 0.25 & Highest weight because first-attempt success is the least retry-inflated signal.\\
\addlinespace
IWSR & 0.20 & Rewards fast convergence and penalizes late retry success.\\
\addlinespace
Exec & 0.20 & Measures structural plan quality even when the goal is not reached.\\
\addlinespace
1 - Hallucination & 0.20 & Captures grounding in the legal action vocabulary.\\
\addlinespace
PAS & 0.15 & Distinguishes sequencing errors from state fabrications.\\
\addlinespace
CoT bonus & +0.05 & Rewards faithful reasoning only when strict CoT alignment exists.\\
\bottomrule
\end{tabularx}
\caption{Composite Planning Score components.}
\label{tab:ps-components-compact}
\end{table}

PS should always be reported with the component table or radar view: the same score can come from very different behaviours, such as strong zero-shot success with weak execution diagnostics or balanced medium performance across all axes.